% 第1章 绪论
\section{绪论}

\subsection{研究背景与意义}

近年来，具身智能（Embodied AI）受到学术界和工业界的广泛关注\cite{duan2022embodied}。与传统的"认知智能"不同，具身智能强调智能体与物理环境的实时交互能力，要求机器人能够通过感知、决策和执行的闭环，在真实世界中完成复杂任务\cite{billard2019learning,kroemer2021review}。这一领域已被纳入国家"十五五"规划建议，成为需要着力打造的新经济增长点，在工业制造、家庭服务、医疗康复、仓储物流等领域拥有巨大的潜在市场。

具身智能的进步同样遵循数据驱动的基本规律。观察认知智能领域的发展历程可以发现，模型性能的提升与训练数据规模呈正相关关系——当语言模型的训练数据从1T增长到十余T时，其理解与生成能力实现了质的飞跃。类似地，机器人学习系统的性能也高度依赖于交互数据的规模与多样性。公开文献显示，早期机器人模型使用数十万条轨迹训练，而近期研究通过引入互联网图文数据实现了更好的泛化能力，也有项目通过百万级异构数据训练获得了通用操作能力。这些数据密集型研究共同指向一个结论：具身智能系统的竞争力取决于数据获取的规模化和多样性保障。

然而，构建大规模的具身数据集面临着现实约束。传统的数据获取方式主要依赖人工遥操作机器人完成任务，这种方案存在多重局限：当采用主从机械臂架构时，需要配置与执行端同构的操控端设备，导致硬件投入成本高昂，且采集场景受限于机械臂的工作空间，此外操作者观察视角与机器人传感器视角的差异可能引入数据偏差；当采用虚拟现实设备遥操作时，虽然降低了对物理机械臂的依赖，但操作者缺乏真实的力触觉反馈，且系统存在网络传输延迟和位姿估计误差，这些因素都会降低操作精度。综合来看，基于机器人本体的传统采集方式在成本、场景适应性和扩展性方面难以满足规模化数据生产的需求。

为突破上述局限，学术界和工业界开始探索不依赖机器人本体的数据获取新范式。这类方法的核心思想是：数据采集者仅通过穿戴或手持专门的传感设备，在真实环境中完成操作演示并记录多模态数据，后续将这些演示数据用于训练控制策略。由于无需机器人本体参与采集过程，这种方法显著降低了数据生产的硬件门槛，使得规模化数据获取成为可能。近年来，多个研究团队推出了相应的采集系统，这些系统通常集成了手持操作端、可穿戴相机和惯性测量单元等组件，形成了轻量化的数据采集解决方案。

本研究旨在构建一套完整的无本体双臂遥操作采集系统，覆盖从数据采集、处理、训练到部署的全流程，并通过跨平台对比实验验证其有效性。研究的主要意义在于：为具身智能研究提供一种低成本、可扩展的数据获取路径，降低数据采集的硬件门槛，推动更大规模具身智能模型的训练与应用落地。

\subsection{研究内容与目标}

本研究围绕无本体数据采集与模仿学习展开，主要研究内容包括以下几个方面：

\textbf{（1）无本体采集系统构建}

设计并实现一套完整的无本体双臂数据采集系统，包括硬件平台选型、软件系统实现、数据格式转换与同步机制。硬件方面采用FastUMI/XV双臂采集设备配合固定位置D435第三视角相机；软件方面实现rosbag到mp4/json再到LeRobot格式的完整数据流程。

\textbf{（2）双臂协调策略研究}

针对双臂协调的挑战，研究有效的数据采集与策略学习方法。探索纯图像输入方案在双臂任务中的可行性，分析放弃显式状态表示（state input）的原因，并以handover任务作为双臂协调的关键验证场景。

\textbf{（3）数据质量分析方法论}

建立数据质量评估体系，设计10条CHECK规则用于数据筛选，实现轨迹相似性分析（基于DTW距离和层次聚类）和图像-动作依赖性分析（基于HSIC互信息估计），为数据采集提供质量监控工具。

\textbf{（4）跨平台实验验证}

通过多组对比实验验证系统有效性：单臂端到端验证（分析数据量-性能关系）、跨平台对比（FastUMI vs XV vs Franka遥操）、第三视角效果分析、双臂handover探索性实验。

本研究的目标是构建一套完整的无本体双臂采集-训练-部署系统，验证其在pick-place等单臂任务和handover等双臂任务上的可行性，为低成本具身数据采集提供可复用的技术方案。

\subsection{论文组织结构}

本文共分为七章，各章内容安排如下：

第一章为绪论，介绍研究背景、意义、内容与目标，以及论文的组织结构。

第二章为相关工作，综述无本体采集方法、双臂机器人遥操作技术、模仿学习策略以及数据质量评估方法的研究现状。

第三章为硬件系统与数据采集，解决"长程任务需求与采集设备头部视角缺失的矛盾"，详细介绍第三视角相机选型（D435固定方案与DJI Nano头戴方案）、多传感器数据同步（含创新的双音chirp标定方案），以及系统实现与性能指标。

第四章为数据处理与质量验证，解决"噪声数据的干扰问题"，阐述数据预处理关键技术（格式转换、降采样、时间对齐）、数据筛选机制，以及轨迹相似性和图像-动作依赖性分析的质量验证方法。

第五章为模型训练与信息嵌入，解决"无本体数据如何实现信息嵌入"，探讨输入设计（纯图像方案）、状态表示探索（ArUco码方案尝试）、核心创新点——Action表示选择（Relative Action vs Absolute Action），以及训练部署流程。

第六章为实验验证与Pipeline体现，展示完整pipeline在单臂任务（pick-place，数据量-性能关系、跨平台对比）和双臂任务（handover，失败模式分析）上的验证结果。

第七章为结论与展望，总结全文工作，指出主要贡献与创新点，并展望后续研究方向。
